\section{Numerical stability of the SDE solver: derivations}\label{app:num-stab}

This appendix supplies the algebra underlying the stability claims
of Section~\ref{sec:num-stab}. We derive the mean-square stability
condition for explicit Euler--Maruyama
(Section~\ref{app:num-stab-test}), apply it to FSA-v2 component by
component (Section~\ref{app:num-stab-fsa}), analyse the role of the
inner substepping for the cubic Stuart--Landau term
(Section~\ref{app:num-stab-substep}), and discuss the interaction
of the discretised CIR diffusion with the Feller boundary condition
(Section~\ref{app:num-stab-feller}).

\subsection{The linear test SDE and explicit Euler--Maruyama}\label{app:num-stab-test}

Consider the linear scalar test SDE
\begin{equation}
  \dd X_t = a\,X_t\,\dd t + b\,X_t\,\dd W_t,
  \qquad a,b \in \R, \quad X_0 \in \R,
  \label{eq:test-sde}
\end{equation}
whose continuous-time solution is the geometric Brownian motion
$X_t = X_0 \exp\!\bigl((a - b^2/2)t + b W_t\bigr)$. Mean-square
stability of \eqref{eq:test-sde} -- i.e.\ $\E[X_t^2] \to 0$ as
$t \to \infty$ -- holds iff
$2a + b^2 < 0$ (Higham, 2000).

The explicit Euler--Maruyama discretisation of \eqref{eq:test-sde}
with step $h$ is
\begin{equation}
  X_{n+1} \;=\; X_n + a\,X_n\,h + b\,X_n\,\sqrt{h}\,\xi_n,
  \qquad \xi_n \sim \Ncal(0,1)\;\text{i.i.d.}
  \label{eq:em-test}
\end{equation}
Squaring and taking expectations,
\begin{align*}
  \E[X_{n+1}^2]
  &= \E\!\bigl[(1 + a h)^2 X_n^2 + 2(1 + a h) b\sqrt{h} X_n^2 \xi_n + b^2 h X_n^2 \xi_n^2\bigr] \\
  &= \bigl((1 + a h)^2 + b^2 h\bigr)\,\E[X_n^2],
\end{align*}
so $\E[X_n^2]$ remains bounded iff
\begin{equation}
  (1 + a h)^2 + b^2 h \;\le\; 1.
  \label{eq:em-msstab-app}
\end{equation}
This is the inequality \eqref{eq:em-msstab} of the main text.
Expanding $(1 + ah)^2 = 1 + 2ah + a^2 h^2$ and rearranging,
\eqref{eq:em-msstab-app} is equivalent to
\begin{equation}
  2a + b^2 + a^2\,h \;\le\; 0,
  \label{eq:em-msstab-rearranged}
\end{equation}
which collapses to the continuous-time condition
$2a + b^2 < 0$ as $h \to 0$. For $a < 0$, the discretisation
imposes the additional constraint
$h \le -(2a + b^2)/a^2$, i.e.\ $h \le 2|a|/(a^2)$ when $b = 0$.

\subsection{Component-by-component bounds for FSA-v2}\label{app:num-stab-fsa}

We linearise the FSA-v2 SDE component by component around the
$A=0$ deterministic equilibrium $(B^*_0,F^*_0,0)$ of
Section~\ref{sec:lyap-branches}. Write $\delta B = B - B^*_0$ etc. The
linearised $B$-equation reads
\begin{equation*}
  \dd(\delta B) \;=\; -(\delta B/\tau_B)\,\dd t + \sigma_B\sqrt{B^*_0(1-B^*_0)}\,\dd W_B
   + \mathrm{(coupling\ terms)}.
\end{equation*}
Treating the coupling terms as forcing, the homogeneous part is the
test SDE \eqref{eq:test-sde} with $a = -1/\tau_B$ and effective
$b^2_{\mathrm{eff}} = \sigma_B^2 B^*_0(1-B^*_0)/(\delta B)^2$. Since
$\delta B$ ranges over the order-unity scale on which we want the
filter to operate, $b_{\mathrm{eff}}^2$ is bounded above by
$\sigma_B^2/4$ (the maximum of $B(1-B)$ on $[0,1]$). Plugging into
\eqref{eq:em-msstab-app},
\[
  (1 - h/\tau_B)^2 + (\sigma_B^2/4)\,h \;\le\; 1
  \;\;\Longleftrightarrow\;\;
  h \;\le\; \frac{2/\tau_B - \sigma_B^2/4}{1/\tau_B^2}
  \;\approx\; 2\tau_B
  \;=\; 84\,\mathrm{d},
\]
where the diffusion contribution $\sigma_B^2/4 \approx 2.5 \times 10^{-5}$
is negligible compared to $2/\tau_B \approx 0.048\,\mathrm{d}^{-1}$.

The linearised $F$-equation has $a = -1/\tau_F^{\mathrm{old}}$
and $b_{\mathrm{eff}}^2 = \sigma_F^2 F^*_0/(\delta F)^2$, bounded
similarly. The same calculation gives
$h \le 2\tau_F^{\mathrm{old}} = 14\,\mathrm{d}$.

The $A$-equation around $A=0$ has $a = \mu(B^*_0,F^*_0)$,
which can be \emph{positive} below the bifurcation locus. In that
case the $A=0$ branch of the deterministic skeleton is unstable
(Proposition~\ref{prop:stab-zero}) and the natural linearisation
point is the oscillatory branch instead, where $a$ becomes
$-2\eta(A^*)^2$. At the limit-cycle amplitude
$A^* = \sqrt{\mu/\eta} \approx 0.81$ (using
$\mu \approx 0.13$, $\eta = 0.20$), we have
$a = -2 \cdot 0.20 \cdot 0.66 \approx -0.27\,\mathrm{d}^{-1}$, and
$b_{\mathrm{eff}}^2 = \sigma_A^2 A^*/(\delta A)^2 \le \sigma_A^2/(A^*)$,
which is also negligible. Plugging in,
$h \le 2/0.27 \approx 7.4\,\mathrm{d}$.

The binding component is therefore the $A$-mode at the limit cycle,
giving the practical stability ceiling of roughly seven days. The
current outer step $\Delta t = 1/96\,\mathrm{d}$ is below this ceiling
by a factor of $700$ -- the quoted ``three orders of magnitude'' of
the main text.

\subsection{The role of the inner substepping}\label{app:num-stab-substep}

The inner-loop substepping of \eqref{eq:em-substep} runs
$n_{\mathrm{sub}} = 4$ deterministic Euler steps of size
$\Delta t / n_{\mathrm{sub}} = \Delta t / 4$ before applying the
single Wiener increment. We now check that the substepping is
neither necessary for stability (which it isn't, given the
factor-700 margin above) nor harmful for it.

A pure deterministic Euler step on the cubic Stuart--Landau ODE
$\dot A = \mu A - \eta A^3$ around the limit-cycle amplitude
$A^* = \sqrt{\mu/\eta}$ has linearised rate
$a_{\mathrm{det}} = -2\eta(A^*)^2 = -2\mu \approx -0.27\,\mathrm{d}^{-1}$.
The deterministic Euler stability condition $|1 + a_{\mathrm{det}} h| \le 1$
becomes $h \le 2/|a_{\mathrm{det}}| \approx 7.4\,\mathrm{d}$, identical
to the bound computed above. For the inner step $h_{\mathrm{sub}} = \Delta t / 4$
to violate this bound would require $\Delta t \ge 30\,\mathrm{d}$, far
beyond any contemplated step.

The substepping is therefore a \emph{accuracy} refinement rather
than a stability device: at large excursions $A \gg A^*$ where the
cubic damping becomes very fast, the inner loop reduces the
deterministic-step error from $O(\Delta t)$ to
$O(\Delta t / n_{\mathrm{sub}})$ at the cost of evaluating the drift
$n_{\mathrm{sub}}$ times per outer step. The Wiener increment is
applied once at the outer-step boundary, so the stochastic
discretisation continues to inherit the strong order-$1/2$
convergence of explicit Euler--Maruyama.

\subsection{The CIR diffusion and the Feller condition}\label{app:num-stab-feller}

The $F$- and $A$-components of the SDE have square-root (CIR-type)
diffusions $\sigma_F\sqrt F\,\dd W_F$ and $\sigma_A\sqrt A\,\dd W_A$.
A continuous-time CIR process
$\dd X = (\alpha - \beta X)\,\dd t + \nu\sqrt X\,\dd W$ remains
strictly positive almost surely if and only if it satisfies the
\emph{Feller condition}
\begin{equation}
  2\alpha \;>\; \nu^2.
  \label{eq:feller}
\end{equation}
For the $F$-component at constant $\Phi = 1$,
$\alpha = \kappa_F \cdot 1 = 0.030$ and $\nu^2 = \sigma_F^2 = 1.44 \times 10^{-4}$,
so $2\alpha/\nu^2 \approx 4.2 \times 10^2 \gg 1$ -- comfortably
satisfied. The $A$-component does not directly fit the CIR template
(the drift is multiplicative, $\mu A - \eta A^3$, not affine), but
the analogous condition near $A = A^*_{\mathrm{osc}}$ involves the
ratio $2\mu A^*_{\mathrm{osc}}/\sigma_A^2 = 2\mu^{3/2}/(\eta^{1/2}\sigma_A^2)$,
also large for the truth values.

The discretised explicit Euler--Maruyama scheme can nevertheless
produce negative values of $F$ or $A$ on rare excursions, even when
\eqref{eq:feller} holds, because of the
$\sqrt{\Delta t}\,\xi_n$ Wiener increment: a single
$|\xi_n| \gg 1$ realisation can push the post-step state below
zero. The reflection step in \eqref{eq:em-substep} catches these
excursions and folds them back into the physiological domain. In
practice the diffusion vanishing at the boundary
($\sigma_F\sqrt F \to 0$ as $F \to 0$, similarly for $A$) makes
such excursions extraordinarily rare; the Feller-aware schemes of
e.g.\ Higham, Mao and Stuart (2003) would handle them more cleanly
but are not currently implemented.

\subsection{What an enlarged step would buy}\label{app:num-stab-cost}

Stage~D and Stage~E5 of Section~\ref{sec:results} ran on the
$15$-minute outer step. The wall-clock cost of those runs was
$\sim\!55$ minutes (Stage~D, $T = 84$ days) and $\sim\!41$ minutes
(Stage~E5, $T = 14$ days, $27$ rolling windows) on a single GPU.
For a fixed simulation horizon $T$, the SDE solver cost scales
linearly in the number of outer steps $T/\Delta t$. Doubling
$\Delta t$ from $15$ to $30$ minutes halves that count and hence
the cost. The full battery of horizon-sweep experiments
(Section~\ref{sec:discussion-other}) at $T \in \{28,42,56,84\}$
days would scale similarly, taking the long $T = 84$-day
closed-loop run from $\sim\!15$ hours down to $\sim\!7.5$ hours
at $\Delta t = 30\,\mathrm{min}$.

The numerical-stability analysis of this appendix supports such an
enlargement up to $\Delta t \approx 1\,\mathrm{hour}$ without a
solver upgrade, conditional on a one-off drift-parity verification
against the $15$-minute reference. Beyond one hour the
$\Phi$-burst-resolution constraint of
Section~\ref{sec:num-stab-binding} starts to dominate, and beyond
that the Lamperti / Milstein / implicit upgrades flagged in
Remark~\ref{rem:num-stab-improvements} become the right route.
